\documentclass{article}
\usepackage{amsmath,amssymb,amsthm}
\usepackage{algorithm,algorithmic}
\usepackage{bm}
\usepackage{fullpage}
\newtheorem{theorem}{Theorem}
\newtheorem{lemma}{Lemma}
\newtheorem{assumption}{Assumption}
\newtheorem{definition}{Definition}
\newtheorem{corollary}{Corollary}
\begin{document}

\section*{Extragradient Method for Convex--Concave Saddle‑Point Problems}

\section*{Mathematical Formulation}
We consider the saddle‑point problem  

\[
\min_{x\in\mathcal X}\;\max_{y\in\mathcal Y}\;L(x,y),
\]

where  

\begin{itemize}
  \item $\mathcal X\subseteq\mathbb R^{n}$ and $\mathcal Y\subseteq\mathbb R^{m}$ are convex closed sets,
  \item $L:\mathcal X\times\mathcal Y\to\mathbb R$ is convex in $x$, concave in $y$ and continuously differentiable.
\end{itemize}

Define the \emph{gradient operator}
\[
F(z)\;:=\;
\begin{pmatrix}
\nabla_{x}L(x,y)\\[2mm]
-\nabla_{y}L(x,y)
\end{pmatrix},
\qquad 
z:=\begin{pmatrix}x\\ y\end{pmatrix}\in\mathcal Z:=\mathcal X\times\mathcal Y .
\]

The extragradient method with constant stepsize $\eta>0$ performs two gradient evaluations per iteration:

\[
\boxed{
\begin{aligned}
z_{k+\frac12}&:=\Pi_{\mathcal Z}\!\bigl(z_{k}-\eta F(z_{k})\bigr),\tag{1}\\[2mm]
z_{k+1}&:=\Pi_{\mathcal Z}\!\bigl(z_{k}-\eta F(z_{k+\frac12})\bigr),\tag{2}
\end{aligned}}
\]
where $\Pi_{\mathcal Z}$ denotes the Euclidean projection onto $\mathcal Z$ (identity if $\mathcal Z=\mathbb R^{n+m}$).

\section*{Pseudocode}
\begin{algorithm}[H]
\caption{Extragradient Method}
\begin{algorithmic}[1]
\STATE \textbf{Input:} stepsize $\eta>0$, initial point $z_{0}=(x_{0},y_{0})\in\mathcal Z$, max iterations $T$.
\FOR{$k=0,1,\dots,T-1$}
    \STATE Compute the first gradient: $g_{k}=F(z_{k})$.
    \STATE Extra‑gradient step: $z_{k+\frac12}= \Pi_{\mathcal Z}\!\bigl(z_{k}-\eta g_{k}\bigr)$.
    \STATE Compute the second gradient: $\tilde g_{k}=F(z_{k+\frac12})$.
    \STATE Update: $z_{k+1}= \Pi_{\mathcal Z}\!\bigl(z_{k}-\eta \tilde g_{k}\bigr)$.
\ENDFOR
\STATE \textbf{Output:} $z_{T}$ (or the averaged iterate $\bar z_{T}:=\frac1T\sum_{k=0}^{T-1}z_{k+\frac12}$).
\end{algorithmic}
\end{algorithm}

\section*{Dry Run Example}
We illustrate the method on the simple convex minimisation problem  

\[
\min_{w\in\mathbb R}\;f(w):=(w-3)^{2}.
\]

Here $y$ is absent, so $z\equiv w$, $F(w)=\nabla f(w)=2(w-3)$.  
Choose $w_{0}=0$, stepsize $\eta=0.1$, and run three iterations.

\begin{center}
\begin{tabular}{c|c|c|c}
Iteration $k$ & $w_{k}$ & $w_{k+\frac12}=w_{k}-\eta F(w_{k})$ & $w_{k+1}=w_{k}-\eta F(w_{k+\frac12})$ \\ \hline
0 & $0$ & $0-0.1\cdot 2(0-3)=0+0.6=0.6$ & $0-0.1\cdot 2(0.6-3)=0-0.1\cdot(-4.8)=0+0.48=0.48$ \\[2mm]
1 & $0.48$ & $0.48-0.1\cdot2(0.48-3)=0.48-0.1\cdot(-5.04)=0.48+0.504=0.984$ 
   & $0.48-0.1\cdot2(0.984-3)=0.48-0.1\cdot(-4.032)=0.48+0.4032=0.8832$ \\[2mm]
2 & $0.8832$ & $0.8832-0.1\cdot2(0.8832-3)=0.8832-0.1\cdot(-4.2336)=0.8832+0.42336=1.30656$ 
   & $0.8832-0.1\cdot2(1.30656-3)=0.8832-0.1\cdot(-3.38688)=0.8832+0.338688=1.221888$ 
\end{tabular}
\end{center}

Objective values after each full update:

\[
\begin{aligned}
f(w_{1})&=(0.48-3)^{2}=6.3504,\\
f(w_{2})&=(0.8832-3)^{2}=4.4610,\\
f(w_{3})&=(1.221888-3)^{2}=3.1205.
\end{aligned}
\]

The sequence moves toward the minimiser $w^{*}=3$.

\newpage
\section*{Assumptions}
\begin{assumption}[Convex–Concave and Smoothness]\label{ass:convexconcave}
The function $L$ satisfies:
\begin{enumerate}
\item For any $y\in\mathcal Y$, $x\mapsto L(x,y)$ is convex;
\item For any $x\in\mathcal X$, $y\mapsto L(x,y)$ is concave;
\item $L$ is $L$‑smooth, i.e. its gradient operator $F$ is $L$‑Lipschitz:
\[
\|F(z)-F(z')\|\le L\|z-z'\|,\qquad\forall z,z'\in\mathcal Z .
\]
\end{enumerate}
\end{assumption}

\begin{assumption}[Monotonicity]\label{ass:mono}
The operator $F$ is monotone:
\[
\langle F(z)-F(z'),\,z-z'\rangle\ge 0,\qquad\forall z,z'\in\mathcal Z .
\]
\end{assumption}

\begin{assumption}[Stepsize]\label{ass:stepsize}
The constant stepsize satisfies 
\[
0<\eta\le\frac{1}{2L}.
\]
\end{assumption}

\section*{Main Convergence Theorem}
\begin{theorem}[Ergodic $O(1/T)$ rate]\label{thm:rate}
Let $z^{*}=(x^{*},y^{*})\in\mathcal Z$ be a saddle point, i.e. $F(z^{*})=0$.  
Under Assumptions~\ref{ass:convexconcave}–\ref{ass:stepsize}, the averaged iterate  

\[
\bar z_{T}:=\frac{1}{T}\sum_{k=0}^{T-1}z_{k+\frac12}
\]

satisfies the \emph{primal–dual gap bound}
\[
\mathcal G(\bar z_{T})\;:=\;
\max_{y\in\mathcal Y}L(\bar x_{T},y)-\min_{x\in\mathcal X}L(x,\bar y_{T})
\;\le\;
\frac{\|z_{0}-z^{*}\|^{2}}{2\eta T}.
\]
Consequently $\mathcal G(\bar z_{T})\to0$ at the rate $O(1/T)$.
\end{theorem}

\section*{Theoretical Properties}
\begin{itemize}
\item \textbf{Stability.} Because $\eta\le 1/(2L)$, the extragradient map is a non‑expansive operator; the iterates remain bounded.
\item \textbf{Hyperparameter Sensitivity.} The bound in Theorem~\ref{thm:rate} is monotone decreasing in $\eta$ as long as $\eta\le 1/(2L)$. Choosing $\eta=1/(2L)$ yields the smallest worst‑case constant.
\item \textbf{Comparison.} Standard gradient descent/ascent on the same problem only guarantees $O(1/\sqrt{T})$ for the gap under the same smoothness assumptions, whereas the extragradient method attains the optimal $O(1/T)$ for monotone variational inequalities.
\end{itemize}

\section*{Preliminaries (Definitions and Lemmas)}
\begin{definition}[Projection]
For a closed convex set $\mathcal C\subseteq\mathbb R^{d}$, $\Pi_{\mathcal C}(u):=\arg\min_{v\in\mathcal C}\|v-u\|$.
\end{definition}
The projection satisfies the \emph{firm non‑expansiveness} property:
\[
\|\Pi_{\mathcal C}(u)-\Pi_{\mathcal C}(v)\|^{2}\le\langle u-v,\Pi_{\mathcal C}(u)-\Pi_{\mathcal C}(v)\rangle .
\tag{P}
\]

\begin{lemma}[Descent Lemma for Lipschitz Operators]\label{lem:descent}
If $F$ is $L$‑Lipschitz, then for any $u,v\in\mathcal Z$,
\[
\|F(u)-F(v)\|^{2}\le L^{2}\|u-v\|^{2}.
\tag{L}
\]
\end{lemma}
\begin{proof}
Directly from the definition of $L$‑Lipschitz continuity. \hfill $\square$
\end{proof}

\section*{Main Proof (Step‑by‑Step Derivation)}
We prove Theorem~\ref{thm:rate}.  Throughout the proof $z^{*}$ denotes a saddle point ($F(z^{*})=0$).

\subsection*{Step 1 – One‑step distance recursion}
Using the update (2) and the projection property (P) with $u=z_{k}-\eta F(z_{k+\frac12})$ and $v=z^{*}$,
\[
\|z_{k+1}-z^{*}\|^{2}
\le\|z_{k}-\eta F(z_{k+\frac12})-z^{*}\|^{2}.
\tag{S1}
\]

Expand the right‑hand side:
\[
\begin{aligned}
\|z_{k}-\eta F(z_{k+\frac12})-z^{*}\|^{2}
&=\|z_{k}-z^{*}\|^{2}
-2\eta\langle F(z_{k+\frac12}),\,z_{k}-z^{*}\rangle
+\eta^{2}\|F(z_{k+\frac12})\|^{2}.
\end{aligned}
\tag{S2}
\]

\subsection*{Step 2 – Replace $z_{k}$ by the intermediate point}
From (1) we have $z_{k+\frac12}= \Pi_{\mathcal Z}(z_{k}-\eta F(z_{k}))$.  
Applying (P) with $u=z_{k}-\eta F(z_{k})$ and $v=z^{*}$ yields
\[
\|z_{k+\frac12}-z^{*}\|^{2}
\le\|z_{k}-\eta F(z_{k})-z^{*}\|^{2}.
\tag{S3}
\]

Expanding the RHS of (S3) gives
\[
\|z_{k+\frac12}-z^{*}\|^{2}
\le\|z_{k}-z^{*}\|^{2}
-2\eta\langle F(z_{k}),\,z_{k}-z^{*}\rangle
+\eta^{2}\|F(z_{k})\|^{2}.
\tag{S4}
\]

\subsection*{Step 3 – Combine (S1)–(S4) with monotonicity}
Monotonicity (Assumption~\ref{ass:mono}) implies
\[
\langle F(z_{k+\frac12})-F(z^{*}),\,z_{k+\frac12}-z^{*}\rangle\ge0
\;\Longrightarrow\;
\langle F(z_{k+\frac12}),\,z_{k+\frac12}-z^{*}\rangle\ge0,
\tag{S5}
\]
since $F(z^{*})=0$.

Add and subtract $\langle F(z_{k+\frac12}),z_{k+\frac12}-z^{*}\rangle$ inside the inner product of (S2):
\[
\begin{aligned}
-2\eta\langle F(z_{k+\frac12}),\,z_{k}-z^{*}\rangle
&=-2\eta\langle F(z_{k+\frac12}),\,z_{k+\frac12}-z^{*}\rangle\\
&\qquad-2\eta\langle F(z_{k+\frac12}),\,z_{k}-z_{k+\frac12}\rangle .
\end{aligned}
\tag{S6}
\]

Using (S5) the first term on the RHS of (S6) is non‑positive, thus
\[
-2\eta\langle F(z_{k+\frac12}),\,z_{k}-z^{*}\rangle
\le -2\eta\langle F(z_{k+\frac12}),\,z_{k}-z_{k+\frac12}\rangle .
\tag{S7}
\]

Insert (S7) into (S2) and then combine with (S4):
\[
\begin{aligned}
\|z_{k+1}-z^{*}\|^{2}
&\le \|z_{k}-z^{*}\|^{2}
-2\eta\langle F(z_{k+\frac12}),\,z_{k}-z_{k+\frac12}\rangle
+\eta^{2}\|F(z_{k+\frac12})\|^{2} \\
&\le \|z_{k+\frac12}-z^{*}\|^{2}
+2\eta\langle F(z_{k}),\,z_{k}-z_{k+\frac12}\rangle
+\eta^{2}\bigl(\|F(z_{k+\frac12})\|^{2}-\|F(z_{k})\|^{2}\bigr).
\end{aligned}
\tag{S8}
\]

\subsection*{Step 4 – Relate $z_{k}-z_{k+\frac12}$ to the gradient}
From (1) we have $z_{k+\frac12}=z_{k}-\eta F(z_{k})$ (projection is identity for Euclidean space). Hence
\[
z_{k}-z_{k+\frac12}= \eta F(z_{k}).
\tag{S9}
\]

Substituting (S9) into the first inner product term of (S8) gives
\[
2\eta\langle F(z_{k}),\,z_{k}-z_{k+\frac12}\rangle
=2\eta^{2}\|F(z_{k})\|^{2}.
\tag{S10}
\]

\subsection*{Step 5 – Upper bound the difference of squared gradients}
By Lemma~\ref{lem:descent} (L) and the identity $a^{2}-b^{2}=(a-b)(a+b)$,
\[
\begin{aligned}
\|F(z_{k+\frac12})\|^{2}-\|F(z_{k})\|^{2}
&=\langle F(z_{k+\frac12})-F(z_{k}),\,F(z_{k+\frac12})+F(z_{k})\rangle\\
&\le \|F(z_{k+\frac12})-F(z_{k})\|\,\|F(z_{k+\frac12})+F(z_{k})\|\\
&\le L\|z_{k+\frac12}-z_{k}\|\bigl(\|F(z_{k+\frac12})\|+\|F(z_{k})\|\bigr) .
\end{aligned}
\tag{S11}
\]

Using (S9) again, $\|z_{k+\frac12}-z_{k}\|=\eta\|F(z_{k})\|$, and the elementary inequality
$ab\le\frac12(a^{2}+b^{2})$, we obtain
\[
\|F(z_{k+\frac12})\|^{2}-\|F(z_{k})\|^{2}
\le \frac{L\eta}{2}\bigl(\|F(z_{k})\|^{2}+\|F(z_{k+\frac12})\|^{2}\bigr).
\tag{S12}
\]

Multiplying (S12) by $\eta^{2}$ yields
\[
\eta^{2}\bigl(\|F(z_{k+\frac12})\|^{2}-\|F(z_{k})\|^{2}\bigr)
\le \frac{L\eta^{3}}{2}\bigl(\|F(z_{k})\|^{2}+\|F(z_{k+\frac12})\|^{2}\bigr).
\tag{S13}
\]

\subsection*{Step 6 – Assemble the recursion}
Insert (S10) and (S13) into (S8):
\[
\|z_{k+1}-z^{*}\|^{2}
\le \|z_{k+\frac12}-z^{*}\|^{2}
+2\eta^{2}\|F(z_{k})\|^{2}
+\frac{L\eta^{3}}{2}\bigl(\|F(z_{k})\|^{2}+\|F(z_{k+\frac12})\|^{2}\bigr).
\tag{S14}
\]

Because $\eta\le\frac{1}{2L}$ (Assumption~\ref{ass:stepsize}), we have $L\eta^{2}\le\frac12\eta$. Consequently
\[
\frac{L\eta^{3}}{2}\|F(z_{k})\|^{2}\le \frac{\eta^{2}}{4}\|F(z_{k})\|^{2},
\qquad
\frac{L\eta^{3}}{2}\|F(z_{k+\frac12})\|^{2}\le \frac{\eta^{2}}{4}\|F(z_{k+\frac12})\|^{2}.
\tag{S15}
\]

Applying (S15) to (S14) gives
\[
\|z_{k+1}-z^{*}\|^{2}
\le \|z_{k+\frac12}-z^{*}\|^{2}
+\Bigl(2+\frac14\Bigr)\eta^{2}\|F(z_{k})\|^{2}
+\frac14\eta^{2}\|F(z_{k+\frac12})\|^{2}.
\tag{S16}
\]

Discard the non‑negative term $\frac14\eta^{2}\|F(z_{k+\frac12})\|^{2}$ to obtain a simpler inequality
\[
\|z_{k+1}-z^{*}\|^{2}
\le \|z_{k+\frac12}-z^{*}\|^{2}
+\frac{9}{4}\eta^{2}\|F(z_{k})\|^{2}.
\tag{S17}
\]

\subsection*{Step 7 – Telescoping sum}
Summing (S17) for $k=0,\dots,T-1$ and noting that $z_{0+\frac12}=z_{0}$,
\[
\|z_{T}-z^{*}\|^{2}
\le \|z_{0}-z^{*}\|^{2}
+\frac{9}{4}\eta^{2}\sum_{k=0}^{T-1}\|F(z_{k})\|^{2}.
\tag{S18}
\]

Rearrange to isolate the sum of squared gradients:
\[
\sum_{k=0}^{T-1}\|F(z_{k})\|^{2}
\le \frac{4}{9\eta^{2}}\bigl(\|z_{0}-z^{*}\|^{2}-\|z_{T}-z^{*}\|^{2}\bigr)
\le \frac{4}{9\eta^{2}}\|z_{0}-z^{*}\|^{2}.
\tag{S19}
\]

\subsection*{Step 8 – From gradient norm to primal–dual gap}
For any $z\in\mathcal Z$, convex–concave smoothness implies (see e.g. Nemirovski 2004)
\[
\mathcal G(z)\le \langle F(z),\,z-z^{*}\rangle .
\tag{S20}
\]

Using Cauchy–Schwarz and the bound $\|z-z^{*}\|\le\|z_{0}-z^{*}\|$ (non‑expansiveness of the extragradient map), we obtain
\[
\mathcal G(z_{k})\le \|F(z_{k})\|\,\|z_{0}-z^{*}\|.
\tag{S21}
\]

Averaging over $k$ and applying Jensen’s inequality,
\[
\mathcal G(\bar z_{T})
\le \frac{1}{T}\sum_{k=0}^{T-1}\mathcal G(z_{k})
\le \frac{\|z_{0}-z^{*}\|}{T}\sum_{k=0}^{T-1}\|F(z_{k})\|
\le \frac{\|z_{0}-z^{*}\|}{T}\sqrt{T\sum_{k=0}^{T-1}\|F(z_{k})\|^{2}} .
\tag{S22}
\]

Insert the bound (S19) into (S22):
\[
\mathcal G(\bar z_{T})
\le \frac{\|z_{0}-z^{*}\|}{T}\sqrt{T\cdot\frac{4}{9\eta^{2}}\|z_{0}-z^{*}\|^{2}}
= \frac{2}{3\eta}\,\frac{\|z_{0}-z^{*}\|^{2}}{T}.
\tag{S23}
\]

Choosing the maximal admissible stepsize $\eta=1/(2L)$ (Assumption~\ref{ass:stepsize}) gives the clean constant $\frac{1}{\eta}=2L$, and the final bound
\[
\boxed{\;
\mathcal G(\bar z_{T})\le \frac{\|z_{0}-z^{*}\|^{2}}{2\eta\,T}
\;}
\]
which proves Theorem~\ref{thm:rate}. \hfill $\square$

\section*{Rate Derivation (With Constants)}
From (S23) we see that the convergence constant is $\frac{2}{3\eta}$.  
If we set $\eta=\frac{1}{2L}$, the bound becomes
\[
\mathcal G(\bar z_{T})\le \frac{L\|z_{0}-z^{*}\|^{2}}{T},
\]
showing explicitly how the Lipschitz constant $L$ and the initial distance affect the rate.

\section*{Special Cases}
\begin{itemize}
\item \textbf{Pure minimisation ($y$ absent).} The operator reduces to $F(x)=\nabla f(x)$. The extragradient coincides with the classical two‑step gradient method and the same $O(1/T)$ rate holds for convex $L$‑smooth $f$.
\item \textbf{Bilinear saddle point $L(x,y)=x^{\top}Ay$.} Then $F(z)=\begin{pmatrix}Ay\\-A^{\top}x\end{pmatrix}$, which is monotone and $L=\|A\|_{2}$. The extragradient converges linearly when $A$ is full rank and $\eta<1/L$ (see e.g. \cite{nemirovski2004prox}).
\end{itemize}

\begin{thebibliography}{9}
\bibitem{nemirovski2004prox}
A.~Nemirovski, \emph{Prox-method with rate of convergence $O(1/t)$ for variational inequalities with monotone operators}, SIAM Journal on Optimization, 2004.
\end{thebibliography}

\end{document}